\subsection{Precision, power and staged validation}
Let $d$ denote the two-arm mean difference in absolute value, divided by the marginal outcome standard deviation, and $d_\psi=\psi/\sigma_Y$ the corresponding standardised matching contrast (half the conventional difference-in-differences). Under equal allocation, independent observations and a common residual variance $\sigma_Y^2(1-R^2)$, where $R^2$ is the variance explained by prespecified pre-assignment covariates, the standardised variance $v$ is $2(1-R^2)/n$ for the two-arm contrast and $(1-R^2)/n$ for the matching contrast, with $n$ retained per cell. The known-variance, two-sided normal-theory power is
\[
\pi(d,v)=\Phi(|d|/\sqrt{v}-z_{1-\alpha/2})
+\Phi(-|d|/\sqrt{v}-z_{1-\alpha/2}).
\]
Here $\Phi$ is the standard normal cumulative distribution function and $z_p$ its $p$ quantile. Recruitment of $K\lceil n/(1-\ell)\rceil$ participants across $K$ cells inflates the target for assumed loss $\ell$; it does not guarantee retention or correct informative missingness.

\begin{table}[htbp]
\centering\small
\caption{Illustrative total recruitment for a single primary contrast: 80\% power, two-sided $\alpha=.05$, 15\% assumed loss.}
\label{tab:power}
\begin{tabular}{@{}rrrrr@{}}
\toprule
&\multicolumn{2}{c}{Two arms}&\multicolumn{2}{c}{Four cells}\\
Effect & $R^2=0$ & $R^2=.25$ & $R^2=0$ & $R^2=.25$\\
\midrule
.20 & 926 & 696 & 928 & 700\\
.30 & 412 & 310 & 416 & 312\\
.40 & 234 & 176 & 236 & 176\\
.50 & 150 & 114 & 152 & 116\\
\bottomrule
\end{tabular}
\end{table}

The companion Python script computes the smallest retained cell size meeting the target and checks the formula by simulating independent Gaussian cell means. At a standardised effect of .30 and $R^2=0$, 175 retained participants per arm and 88 per factorial cell yield simulated powers of .80065 and .80200, respectively, in 20,000 repetitions per design. At the same sample sizes under the null, rejection proportions are .04865 and .04925 (20,000 repetitions each). Monte Carlo standard errors for the power estimates are approximately .0028.

The sensitivity analysis assumes known variance. For a balanced Gaussian design with common estimated variance, the corresponding noncentral-$t$ distribution has $K(n-1)$ degrees of freedom; at $d=.30$ and $R^2=0$, the minimum is 176 per arm or 88 per factorial cell. A supplementary Python diagnostic uses Welch--Satterthwaite intervals with estimated cell variances, unequal allocation, heterogeneous variance and independent retention; it reports bias, rejection, coverage, valid replicate counts and Monte Carlo uncertainty. Separate synthetic scenarios examine clustered diary slopes, expertise-by-interface interactions and enjoyment trends. Extended mixed-likelihood diagnostics (2,000 replicates per scenario; 80 participants, eight observations) yielded H6 null rejection .052 without missingness and .0815 under MNAR. H10 power at an illustrative contrast of $-.35$ was .193 without missingness and .133 under MAR; covariance boundary solutions were frequent. These findings preclude confirmatory H10 at that size. Revised synthetic diagnostics (2,000 replicates per cell) support conservative H6 CR2 inference under MAR (null rejection .037 at $\alpha=.04$), not unrestricted MNAR (.067). With active final-session recovery (approximately 83\% observed under MAR), the revised endpoint estimator at 800 assigned participants yields .868 power for an illustrative contrast of $-.35$; 600 yields .730. This is a conditional planning scenario, not an empirically justified recruitment target. Recalculate power and calibrate inference for the selected estimator, clustering, attrition and multiplicity. Choose the smallest effect of interest from outcome meaning, not pooled effects of other musical activities. Begin with qualitative discovery, cognitive interviews and feasibility work before a focused confirmatory experiment. Pilot data should inform measurement properties, variance and completion rates, not determine the target effect through an unstable preliminary estimate.


\paragraph{Rounded recruitment sensitivity (27 November 2026).}
For the illustrative factorial matching contrast with $d_\psi=.30$, no baseline variance reduction and 15\% independent loss, the proposed $N=420$ allocates 105 participants per cell and yields 357 expected observed outcomes. The noncentral-$t$ approximation gives power .80689. An independent local reproduction generated 100,000 Gaussian cell-summary datasets at each of the null and alternative, with binomial retention, estimated cell variances and Welch--Satterthwaite intervals. Null rejection was .05084; alternative rejection was .80578 (Monte Carlo SE .00125); coverage was .94916 and .94982, respectively. Mean retained $N$ was 356.93 and 357.02. These calculations do not calibrate bounded outcomes, informative loss or the remaining hypotheses. The recruitment target remains illustrative.


\subsection{Joint prediction calibration}
A separate synthetic screen evaluated the intersection-union decision for $s_0>0$ and $s_1>0$. Gaussian independent cell summaries were generated with binomial outcome-independent retention and estimated variances. The grid crosses six operational sample sizes (420, 600, 800, 1000, 1200, 1400), eight symmetric, asymmetric, global-null, partial-null and reversed simple-effect configurations, and two variance/retention profiles (20,000 independent datasets per cell, 1,920,000 total; seed 20261203). The equal profile uses unit residual standard deviations and .85 retention; the heterogeneous profile uses standard deviations (.8, 1.2, 1.2, .8) and retention (.90, .80, .75, .90) in cell order 00, 01, 10, 11. All simulated fits were valid. At $s_0=s_1=.30$ in common unit-scale points, joint rejection probabilities for $N=420,800,1000,1200$ were .40685, .75240, .85760 and .91910 in the equal profile, and .37655, .70515, .81770 and .88940 in the heterogeneous profile. For $N=420$, an asymmetric (.15,.45) configuration had joint power .23660 despite mean-contrast rejection .80730. A partial-null (0,.60) configuration produced joint rejection .04970, whereas mean-contrast rejection was .80910. Across the simulated union-null configurations the largest joint rejection rate was .05225, within Monte Carlo uncertainty of .05; simple-effect pointwise coverage ranged .94525 to .95385 across all cells. At $N=1000$, joint power Monte Carlo standard errors were .00247 and .00273 for the two symmetric profiles. These simulations do not establish that the bounded rating is Gaussian, invariant or missing at random. The symmetric-effect grid identifies $N=1000$ as the first evaluated candidate exceeding .80 in both profiles; it does not establish a minimum for asymmetric effects or a recruitment target. Figure~\ref{fig:cp-power} displays the reported subset of joint-power estimates and contrasts the symmetric, asymmetric and partial-null configurations at $N=420$. The mean-contrast rejection probability is not the power of CP. Preserve the earlier mean-contrast results as a distinct historical estimand. Pilot evidence, justified component effect sizes, measurement comparability, informative-missingness analysis and owner approval remain necessary. The local execution used Python numerical dependencies NumPy 2.3.5 and SciPy 1.17.0; repository-pinned CI reproduction remains separate.

\begin{figure}[htbp]
\centering
\begin{tikzpicture}[x=1cm,y=1cm,font=\small]
\begin{scope}[xscale=0.006,yscale=4]
\draw[->] (380,0)--(1245,0) node[below right] {$N$};
\draw[->] (380,0)--(380,1.05);
\foreach \v in {0,.2,.4,.6,.8,1}{\draw[gray!25] (380,\v)--(1220,\v);\node[left] at (380,\v) {\v};}
\foreach \n in {420,800,1000,1200}{\node[below] at (\n,0) {\n};}
\draw[dashed] (380,.8)--(1220,.8);
\draw[thick,blue!70!black] plot coordinates {(420,.40685)(800,.75240)(1000,.85760)(1200,.91910)};
\draw[thick,orange!85!black,densely dotted] plot coordinates {(420,.37655)(800,.70515)(1000,.81770)(1200,.88940)};
\foreach \n/\v in {420/.40685,800/.75240,1000/.85760,1200/.91910}{\fill[blue!70!black] (\n,\v) circle[radius=1.5pt];}
\foreach \n/\v in {420/.37655,800/.70515,1000/.81770,1200/.88940}{\draw[orange!85!black,fill=white] (\n,\v) circle[radius=1.5pt];}
\node[above,align=center,font=\small\bfseries] at (800,1.05) {A. Joint power};
\node[above right,blue!70!black,font=\scriptsize] at (420,.45) {Equal};
\node[below right,orange!85!black,font=\scriptsize] at (420,.32) {Heterogeneous};
\end{scope}
\begin{scope}[xshift=8.1cm,x=1cm,y=4cm]
\draw[->] (0,0)--(4.7,0);
\draw[->] (0,0)--(0,1.05);
\foreach \v in {0,.2,.4,.6,.8,1}{\draw[gray!25] (0,\v)--(4.5,\v);\node[left] at (0,\v) {\v};}
\fill[blue!70!black] (.65,0) rectangle (1.25,.23660);
\fill[gray!55] (1.28,0) rectangle (1.88,.80730);
\fill[blue!70!black] (2.65,0) rectangle (3.25,.04970);
\fill[gray!55] (3.28,0) rectangle (3.88,.80910);
\node[below,align=center,font=\scriptsize] at (1.25,0) {(.15,.45)};
\node[below,align=center,font=\scriptsize] at (3.25,0) {(0,.60)};
\node[above,font=\small\bfseries] at (2.3,1.05) {B. Rejection at $N=420$};
\node[blue!70!black,font=\scriptsize] at (1.05,.95) {Joint};
\node[font=\scriptsize] at (3.3,.95) {Mean contrast};
\end{scope}
\end{tikzpicture}

\caption{Reported synthetic calibration results. Panel A: joint rejection probabilities for symmetric simple effects (.30, .30) at four reported sample sizes. Lines connect evaluated points only; the dashed line is the .80 planning threshold. Panel B: joint and mean-contrast rejection probabilities at $N=420$ for asymmetric and partial-null configurations. Under the partial null the joint rejection probability is a Type I error rate, not power. No participant observations are represented.}
\label{fig:cp-power}
\end{figure}

\subsection{Bounded-outcome and observation sensitivity}
A supplementary synthetic design diagnostic represents the candidate rating as an integer from 0 to 10 using a beta-binomial distribution. The full-population simple effects are defined on the observed score scale, not a latent Gaussian scale. Low, intermediate and high baseline means assess floor and ceiling behaviour; two concentration parameters vary score dispersion. A binary pre-assignment prognostic factor influences both the score and observation under MAR, whereas MNAR additionally depends on the unobserved score. The diagnostic compares available-case Welch inference with equally weighted stratum means and Satterthwaite uncertainty under a known target stratum distribution. Standardisation addresses the specified MAR mechanism under adequate stratum support, but does not identify MNAR outcomes. The planned grid crosses three sample sizes, five effect patterns, three baseline means, four observation mechanisms and two dispersion settings. Report unconditional conjunction rejection, valid and failed counts, Monte Carlo standard errors and conditional bias, RMSE and coverage. Results are not reported here until exact-source execution is verified; the diagnostic does not establish a validated rating, a recruitment target or an approved confirmatory estimator.
